\documentclass[11pt,a4paper]{article}

\usepackage[spanish]{babel}
\usepackage[utf8]{inputenc}
\usepackage[T1]{fontenc}

\usepackage[a4paper,margin=2.5cm]{geometry}

\usepackage{xcolor}
\usepackage{tcolorbox}
\usepackage{hyperref}
\usepackage{listings}
\usepackage{enumitem}
\usepackage{titlesec}

\hypersetup{
	colorlinks=true,
	linkcolor=blue,
	urlcolor=blue
}

\definecolor{godotblue}{RGB}{71,114,179}
\definecolor{lightblue}{RGB}{235,242,250}
\definecolor{codebg}{RGB}{245,245,245}

\titleformat{\section}
{\Large\bfseries\color{godotblue}}
{\thesection.}{0.5em}{}

\lstset{
	backgroundcolor=\color{codebg},
	basicstyle=\ttfamily\small,
	frame=single,
	breaklines=true,
	showstringspaces=false
}

\begin{document}

\begin{center}

{\Huge\bfseries Guía Rápida}\\[0.4cm]

{\LARGE Recursos Interactivos con Godot}\\[0.3cm]

{\large Exportación Web + GitHub Pages}

\vspace{0.5cm}

\begin{tcolorbox}[
	colback=lightblue,
	colframe=godotblue,
	width=0.9\textwidth
]
Esta guía recoge los pasos mínimos necesarios para publicar una figura
interactiva desarrollada en Godot utilizando GitHub Pages.
Está pensada como recordatorio rápido para futuros proyectos.
\end{tcolorbox}

\end{center}

\section{Exportación desde Godot}

Una vez finalizado el proyecto:

\begin{enumerate}[label=\arabic*.]
	\item Abrir:
	\texttt{Proyecto $\rightarrow$ Exportar}
	
	\item Seleccionar:
	\texttt{Web}
	
	\item Exportar el proyecto.
	
	\item Utilizar como nombre:
	
	\begin{center}
		\Large\texttt{index}
	\end{center}
	
	De este modo Godot generará automáticamente:
	
\begin{lstlisting}
index.html
index.js
index.wasm
index.pck
\end{lstlisting}

	\item Comprobar que todos los archivos se encuentran en la carpeta de exportación.
\end{enumerate}

\section{Creación del repositorio en GitHub}

Crear un repositorio nuevo en GitHub:

\begin{itemize}
	\item Público.
	\item Sin README.
	\item Sin .gitignore.
	\item Sin licencia.
\end{itemize}

Por ejemplo:

\begin{center}
\texttt{fenomeno-runge}
\end{center}

La URL SSH tendrá la forma:

\begin{lstlisting}
git@github.com:eduroyo/fenomeno-runge.git
\end{lstlisting}

\section{Inicialización del repositorio local}

Situarse en la carpeta exportada y ejecutar:

\begin{lstlisting}
git init
git add .
git commit -m "first commit"
git branch -M main
git remote add origin git@github.com:eduroyo/fenomeno-runge.git
git push -u origin main
\end{lstlisting}

\section{Actualizaciones posteriores}

Si el repositorio ya existe y únicamente se desea actualizar el contenido:

\begin{lstlisting}
git add .
git commit -m "actualizacion"
git push
\end{lstlisting}

\section{Activar GitHub Pages}

En el repositorio:

\begin{center}
\texttt{Settings $\rightarrow$ Pages}
\end{center}

Seleccionar:

\begin{itemize}
	\item Source: \texttt{Deploy from a branch}
	\item Branch: \texttt{main}
	\item Folder: \texttt{/ (root)}
\end{itemize}

Guardar los cambios.

\section{URL final}

Tras unos minutos la aplicación quedará accesible en:

\begin{lstlisting}
https://eduroyo.github.io/fenomeno-runge/
\end{lstlisting}

\section{Comprobaciones rápidas}

Verificar conexión SSH:

\begin{lstlisting}
ssh -T git@github.com
\end{lstlisting}

La respuesta debe ser similar a:

\begin{lstlisting}
Hi eduroyo! You've successfully authenticated...
\end{lstlisting}

Consultar la rama actual:

\begin{lstlisting}
git branch
\end{lstlisting}

Consultar el repositorio remoto:

\begin{lstlisting}
git remote -v
\end{lstlisting}

\section{Notas}

\begin{tcolorbox}[
	colback=lightblue,
	colframe=godotblue
]
\begin{itemize}
	\item Utilizar siempre nombres sencillos para los repositorios.
	\item Mantener la rama principal como \texttt{main}.
	\item Comprobar que los archivos \texttt{.wasm} y \texttt{.pck} se han subido correctamente.
	\item Si Git pide usuario y contraseña, probablemente se esté usando HTTPS en lugar de SSH.
	\item La URL remota debería comenzar por:
	
\begin{center}
\texttt{git@github.com:}
\end{center}

\end{itemize}
\end{tcolorbox}

\vfill

\begin{center}
\rule{0.5\textwidth}{0.4pt}

\vspace{0.2cm}

\textit{Programa de Incentivación de la Innovación Docente en la Universidad de Zaragoza}\\
\textit{Matemáticas visuales: recursos interactivos para transformar el aprendizaje en grados STEM}

\end{center}

\end{document}